% Configuración común del informe.
\usepackage[T1]{fontenc}
\usepackage[utf8]{inputenc}
\usepackage[spanish,es-noquoting]{babel}
\usepackage{newtxtext,newtxmath}   % Times New Roman (texto y ecuaciones)
\usepackage{graphicx}
\usepackage{amsmath}
\usepackage{booktabs}
\usepackage{caption}
\usepackage{subcaption}
\usepackage[margin=2.5cm]{geometry}
\usepackage{siunitx}
\usepackage{hyperref}

\hypersetup{hidelinks}
\sisetup{output-decimal-marker={,}, separate-uncertainty=true}

% "Figura 1: ..." y "Tabla 1: ..." con la etiqueta en negrita.
\captionsetup{labelfont=bf, font=small, width=0.9\textwidth}

% Abreviaturas de referencia cruzada usadas en el texto.
\newcommand{\figref}[1]{Fig.~\ref{#1}}
\newcommand{\ecref}[1]{Ec.~(\ref{#1})}
\newcommand{\tabref}[1]{Tabla~\ref{#1}}

% Vectores: negrita sin itálica. Escalares: itálica sin negrita.
\renewcommand{\vec}[1]{\mathbf{#1}}

% Babel en español rotula las tablas como "Cuadro"; el informe usa "Tabla".
\addto\captionsspanish{\renewcommand{\tablename}{Tabla}}
